\chapter{Introduction}
\label{chap:intro}

The combination of existing scientific disciplines provides fertile ground
for the development of entirely new scientific discoveries. Computer science
and quantum physics are now coming together in this way to give birth to
quantum computing. Computer science and quantum physics originated
as contemporaries in the early part of the 20th century, culminating in
technology that decisively ended World War II. They have remained
largely independent until recently, but the common thread connecting them
has always existed.

From the perspective of computer science, engineering is problem-driven.
We humans wish to solve a problem, like calculating ballistic weapon
trajectories, and we work backwards from there. To calculate,
we need to perform fast
arithmetic on stored numbers, for which we need to design devices in a
certain configuration, for which we need certain materials from nature.
At heart, it is the study of structured problem-solving which happens to have a strong
mechanical bias. This perspective asks the questions: what do we want to do,
and how can we
do it? It is like a child going into a LEGO store with a pre-existing plan
and only buying the bricks needed. It is an efficient plan for brick procurement,
but not
necessarily creative.

From the perspective of quantum physics, engineering is phenomena-based.
We humans wish to describe reality and observe nature. We discover objects
and effects in nature, and afterwards we may
then discover that they
can be useful in performing some useful task. An example is
the discovery of quantum tunnelling and its subsequent use to build a
field-effect transistor (FET) for digital electronics. This perspective asks the
questions: what is available to us, and what could we do it? It is like a child
going into a LEGO store with no pre-conceived plan, investigating all the
parts and dreaming up ways to use them. This will generate completely new
designs and models, but it is much slower than brick-shopping.

These perspectives are intertwined and serve complementary purposes,
and we often alternate
between the two of them. In our transistor example, we probably would not
have thought to build a FET if we had not already needed fast electrical
switching, for which vacuum tubes were proving inadequate. Moreover, we can
go beyond using quantum physics to simulate a classical device. There is
a remarkable example
of a completely new algorithm designed using the building blocks of
quantum physics. In 1994, Peter Shor contributed a quantum factoring algorithm
to solve a human problem efficiently,
one that has defied all classical algorithms before or since.
where no known classical algorithm, previously or since.
This kickstarted the entire field of quantum computing; united many
disciplines including computer science, physics, engineering, and
mathematics; and provided the foundation for this dissertation.

The problem of factorization, or factoring,
is simple to state and difficult to solve.
Given an integer $m$, return the integers $\{p_1, \ldots, p_k\}$ whose
product of their powers is $m$: $m = p_1^{\alpha_1}\cdots p_k^{\alpha_k}$.
It is a beautiful theoretical problem posed by Carl Friedrich Gauss in 1801,
one which he felt that every self-respecting scientist should devote their
efforts to solving. Factoring is an important problem in the theory of
numbers, which Gauss considered the crown jewel of mathematics, itself
considered the queen of the sciences.\footnote{\emph{Gauss zum Ged\"achtniss} (1856) by Wolfgang Sartorius von Waltershausen}
It is a quirk of cryptography that the hardness of problems can be used
to secure communication. Therefore, factoring also has immense practical
usefulness as the basis of the RSA cryptosystem, the most widely-used
cryptosystem around the world and on the internet. The somewhat sinister,
but financially important, importance to
governments is the ability to gather intelligence on each other and on their
citizens. The importance to citizens is the ability to choose security
parameters so that their correspondence is protected for a limited time.
The organizers of the popular democratic uprisings in the Middle East known
as the Arab Spring definitely benefitted from using a cryptosystem that was
secure against compromise given the resources of the regimes that they toppled.
Similarly, being able to accurately estimate the resources needed to factor
an RSA key of a certain size may help protect future citizens against
governments with a quantum computers, as governments will be the first
organizations interested in and able to afford such a device.

One of the most remarkable scientific developments of the last twenty years is
 limited not just to one field of human knowledge, but it is exists at the
 confluence of many. Quantum computation, and quantum information, is the
 study of the power and ultimate limits of human abilities in the realm of
 computation and information processing. In particular, it applies to any
 machine that humans can build, given that all machines are physical objects
 and are therefore governed by the laws of physics. One definition of physics
 is the study of reality, or how to explain and predict the behavior of
 phenomenona in our universe. One definition of computer science is the study
 of problem-solving, or how information can be manipulated. From a human
 perspective, we can narrow such definitions: physics studies the phenomena
 which are observable to humans, at least in principle, given a long enough
 time and enough resources; computer science studies those problems which are
 solvable by humans and those problem-solving machines which are buildable
 and runnable by humans, again, given enough time and resources.
 This perspective also makes the union of physics and computer science
 (as well as mathematics, engineering, and other related fields) natural,
 although such a connection was not made until 1994.

The field of quantum computing effectively began that year when Peter Shor discovered that exploiting the mathematics of quantum physics, one could in theory build a quantum computer to determine the integer factors of a large number efficiently. This problem, known as factoring, resisted all attempts at solution from the time it was recognized by Karl Friedrich Gauss as the
problem ``to which we should bend our every effort'', the greatest problem of number theory, the subfield which he considered the jewel of mathematics, itself the queen of all the sciences.

Factoring, in addition to being a beautiful abstract problem, also has a very real implication for human affairs. Because of its difficulty on current digital computers, factoring has become the basis of the widely-used RSA cryptosystem, used to secure most of the worlds online transactions.

So now we have discovered that quantum computing is theoretically interesting because it can solve an important human problem. But can quantum computers be actually built?

Quantum architecture is the intermediate layer between algorithms and hardware.
It is the design of how quantum bits and their allowed interactions in order to solve these algorithms efficiently on realistic models of quantum hardware.
It aims to minimize circuit resources of interest, namely depth, size, and width. In analogy to classical circuits, the depth is the running time of an algorithm allowing parallelization, the size is the number of operations required over all parallel processors, and the width is the number of (quantum) bits required over all parallel processors.

In this dissertation, we study quantum architecture in the context of optimizing Shor's factoring algorithm on nearest-neighbor architectures with realistic constraints. It is hoped that lessons learned in this special-case can contribute to the general community general principles which can be used to generalize other quantum algorithms. We posit that the larger overall goal of quantum architecture should be the design of a general-purpose quantum processor, one which can execute any quantum algorithm with emphasis on being able to perform a core group of operations efficiently. This core group of operations is defined by the instruction set. Once quantum architecture has progressed to this point, we can leverage the remarkable strides in normal digital architecture over the past 80 years.

As well, we may be able to use the insights of quantum architecture to build quantum processors which can surpass digital architectures in the solution of large-scale problems over exponentially-sized solution spaces.

In Chapter \ref{chap:factor-polylog}, we present our first main result, a
nearest-neighbor architecture for factoring in polylogarithmic depth. This is an exponential improvement of the previous best-known result, which required quadratic depth.

In Chapter \ref{chap:factor-sublog}, we further improve our result with a
factoring architecture which executes in sublogarithmic depth, which represents another exponential improvement over our first result. Therefore, this achievement represents a remarkable doubly-exponential improvement over previous factoring architectures.

In Chapter \ref{chap:qcompile}, we study the relationship between quantum architecture and compiling. In digital computing, the boundary between architecture and compilers is quite porous and is determined by a processor's instruction set. Architecture studies processor resources to solve an algorithm given a particular instruction set which is fixed in hardware. This instruction set is produced by a compiler, a piece of (low-level) software which transforms over pieces of (high-level) software. This instruction set can change based on which algorithms it allows to solve efficiently as well as which processors it allows to manufacture efficiently as well as which operations it allows humans to understand easily. All of these factors combine to make architecture an art and an engineering discipline rather than merely a science.

Quantum computers make this problem even more difficult due to the nature of a quantum bit. Because transformations between quantum states vary continuously over the space of unitary matrices with complex coefficients, we can only approximate desired quantum logic gates using a fixed set, given to us by fault-tolerance.

This difficulty of approximation, which we call quantum compiling, represents the fundamental limit to the depth of factoring architectures. All known constant-depth factoring implementations, as well as sorting networks for those implementations which are not nearest-neighbor, assume arbitrary single-qubit rotations. That is, they assume the model where quantum compiling can be done in constant-depth, in polynomial size, for free. However, this is not the case when we combine these abstract circuit models with our model of fault-tolerant quantum computing. Under FTQC, quantum compiling depth represents a sufficient, but not known to be necessary, condition for constant-depth factoring. 

We leave this ultimately as an open problem and change directions.
Given that decreasing factoring depth with abandon increases size and width at relatively large polynomial growth rates (up to $n^8$). This is a rather undesirable time-space tradeoff.

Here, we insert the estimate of electricity and square miles and centuries of run-time based on our estimates given in the final exam.

In Chapter \ref{chap:coherence}, we introduce a novel quantum circuit resource called \emph{circuit coherence}.

% Point 
We are now ready to state the thesis that we investigate in this dissertation.

\begin{quote}
Studying the depth of quantum architectures can help us design quantum computers that can solve
human problems with a human lifetime.
\end{quote}

Here, or somewhere else, we must distinguish between quantum architecture as a field and quantum architecture as a
circuit with hardware constraints. Actually, the two definitions are pretty close.

We also need to provide preliminary definitions of a quantum bit.

This deserves some unpacking. A quantum architecture is a circuit designed to implement a particular algorithm given certain
hardware constraints. The depth of a quantum architecture is the amount of time it takes for the circuit to complete.
Given a particular input $x$, how long (in human clock time) does it take to return the output of some function $f(x)$?
Quantum computers are interesting because for some algorithms, namely factoring, the depth of a quantum circuit is
less than the depth of any known classical circuit. However, this depth may correspond to a running time which is
still much longer than a single human lifetime. In particular this author, who is 32 years old at the time of this
writing, is interested in techniques to optimize the successful completion of a problem within his own lifetime.

Therefore, one meaning of ``studying the depth'' means to \emph{decrease} depth. The main technique for decreasing
depth is to parallelize a quantum circuit, that is, to divide and duplicate information to allow them to be processed
in a parallel way, using multiple classical controllers. However, parallelizing can increase other circuit resources,
usually due to this information duplication. This introduces a time-space tradeoff. In practical reality, we cannot
decrease depth at the expense of other circuit resources, such as the amount of information in a form suitable for
operation over time.

Decreasing depth can increase circuit size (the total number of operations to be performed) and circuit width
(the total amount of quantum information that can exist at any one time). This increase may again far
outstrip human resources, and delay completion of a quantum computer's construction, let alone its
successful execution, until well after that author's death. Therefore, the useful study of depth to help
quantum computers solve human problems must also limit this time-space tradeoff overhead. It should also
involve defining the exact tradeoff that is useful to minimize.

We have outlined the problems of designing such a low-depth quantum architecture so far in a very general,
abstract form. In this thesis, we will make these concepts more concrete by describing a formalism for
quantum circuits and quantum computation that is well-known to the research community.
Building on this formalism, we will propose a new architectural model and present our solutions to factoring in them.


Section \ref{sec:intro-arch} will discuss architectural models. We will present the definition of an
architectural models and give examples of them from the
current literature, with names like \textsf{AC} and \textsf{NTC}. We will discuss the idea of architectural
sub-models, with names like \textsf{2D CCNTC}. Finally,
we will describe our contribution of a new architectural model, called \textsf{2D CCNTCM}, which is a weakly 2D,
hybrid architecture that introduces the idea of allowing both short-range (local) interactions and long-range
interactions, and providing a configurable tradeoff in between them.

In order to build long-range interactions out of short-range interactions, Section \ref{sec:intro-cdc} 
gives a thorough background on recent constant-depth communication techniques. These include three
basic constructs for teleporting, copying, and uncopying quantum information across an unbounded
lattice.

All operations on a quantum architecture can be reduced to single-qubit operations and two-qubit operations.
The latter operations are what we have been calling ``interactions,'' and we have studied how to overcome
constraints on nearest-neighbor interactions using the constant-depth communication techniques previously.

Ultimately, however, the limits to depth of a quantum architecture are dependent on the single-qubit operations.

\section{Quantum Gates and Circuit Bases}
\label{sec:intro-basis}

To compile, or implement, arbitrary quantum algorithms, we must construct circuits
out of gates from a universal set, which we call a \emph{circuit basis},
or just \emph{basis}.
This should not be confused with a basis for a vector space.
Therefore, we will now review quantum gates and how to combine them into
circuits. This procedure is known as quantum compiling, but we will not
delve into its details until we need them in Chapter \ref{chap:qcompile}.

A quantum gate on $n$ qubits is a $2^n \times 2^n$ unitary matrix
(an element of $U(2^n)$). We can consider this the overall circuit width.
Often, we find it useful to neglect a
global phase, since these cannot be measured in quantum mechanics.
However, a global
phase on a particular system $S$ may result in a measurable relative phase
in a large system $S'$ of which $S$ is a subsystem. Therefore, for our
purposes we will only distinguish between $U(2^n)$ and
$SU(2^n)$ in the few cases where it matters for
quantum compiling. The distinction between a quantum circuit
and a quantum gate is relative; often we consider a quantum gate as a
fundamental primitive of our physical technology, and a circuit as a
composite of these gates corresponding to a quantum algorithm.

In Section \ref{subsec:pauli} we
will review the Pauli single-qubit gates and their corresponding group.
In Section \ref{subsec:clifford} we will introduce the Clifford group.
In Section \ref{subsec:controlled} we will introduce controlled operations
and the Toffoli gate.
In Section \ref{subsec:qcompile-single} we will discuss \emph{single-qubit compiling}
and how a
general single-qubit gate can be compiled into rotations about Bloch sphere axes.
In Section \ref{subsec:distance} we will present distance metrics to
measure the quality of our single-qubit (and later multi-qubit) approximations.
In Section \ref{subsec:qcompile-bases} we will finally
define what it means for a gate set to be universal.

%%%%%%%%%%%%%%%%%%%%%%%%%%%%%%%%%%%%%%%%%%%%%%%%%%%%%%%%%%%%%%%%%%%%%%%%%%%%%%
\subsection{Pauli Group}
\label{subsec:pauli}

We review here the Pauli group on one qubit, $\mathcal{P}_1 = \{I, X, Y, Z\}$.
These last three represent
rotations of $\pi$ on the Bloch sphere about the $x$-axis, $y$-axis, and $z$-axis,
using the homomorphism between $SU(2)$ and $SO(3)$. The $2\times 2$ identity matrix
is denoted $I$. The group $\mathcal{P}_1$ also serves as a complex vector
basis for generating elements of $U(2)$.

\begin{IEEEeqnarray}{rcrclCrcrcl}
X & = & \sigma_x & = &
 \left[
  \begin{array}{cc}
    0 & 1 \\
    1 & 0 \\
  \end{array} \right]
& \qquad &
Y & = & \sigma_y & = &
 \left[
  \begin{array}{cc}
    0 & i \\
   -i & 0 \\
  \end{array} \right]
\\
Z & = & \sigma_z & = &
 \left[
  \begin{array}{cc}
    1 & 0 \\
    0 & -1 \\
  \end{array} \right]
& \qquad &
I & = & \sigma_0 & = &
 \left[
  \begin{array}{cc}
    1 & 0 \\
    0 & 1 \\
  \end{array} \right]
\end{IEEEeqnarray}

We define the Pauli group $\mathcal{P}_n$ on $n$ qubits as the set of
all $n$-qubit operators which are tensor products of elements from
$\mathcal{P}_1$.

%%%%%%%%%%%%%%%%%%%%%%%%%%%%%%%%%%%%%%%%%%%%%%%%%%%%%%%%%%%%%%%%%%%%%%%%%%%%%%
\subsection{The Clifford Group}
\label{subsec:clifford}

We define the normalizer of $\mathcal{P}_n$ as the
Clifford group $\mathcal{C}_n$ on $n$ qubits.

\begin{equation}
\mathcal{C}_n = \{ C \in U(2^n) | CPC^{\dagger} \in \mathcal{P}_n \quad \forall P \in \mathcal{P}_n \}
\end{equation}

Of particular interest to us is the two-qubit Clifford group $\mathcal{C}_2$,
which is generated by the following matrices:

\begin{equation}
\mathcal{C}_2 = \langle H, S, CNOT \rangle
\end{equation}

The first two Clifford generator matrices are single-qubit gates ($2 \times 2$ unitary matrices) and
their inclusion means they can be applied on either the first or the second
qubit.\footnote{Standard convention writes $H_i$ to mean $H$ on qubit $i$ and likewise for $S_i$.}
The matrix $H$ is known as the Hadamard gate, and it is a special case of the
general Walsh-Hadamard transform. It is its own adjoint: $H^{\dagger} = H$.
The matrix $S$ is known as the phase gate, and it can be considered the
``square root'' of the Pauli $Z$ gate (up to a phase): $S^2 = Z$.
Equivalently, it can be viewed as a $\pi/2$ rotation about the Bloch sphere
$z$-axis, and its adjoint $S^{\dagger}$ is the reverse rotation of $-\pi /2$.
These matrices are defined below.

\begin{equation}
H = \normtwo
 \left[
  \begin{array}{cc}
    1 & 1 \\
    1 & -1 \\
  \end{array} \right]
\qquad
S = 
 \left[
  \begin{array}{cc}
    1 & 0 \\
    0 & i \\
  \end{array} \right]
\qquad
S^{\dagger} = 
 \left[
  \begin{array}{cc}
    1 & 0 \\
    0 & -i \\
  \end{array} \right]
\end{equation}

The Hadamard matrix also has the special property that it changes between the
$X$ basis and the $Z$ basis, that is, the vector basis for single-qubit
states consisting of eigenstates of the Pauli $X$ and Pauli $Z$ gates,
respectively. In fact, using the identities $X = HZH$ and $S^2 = Z$, it
is easy to see why $X$ and $Z$ are often listed as generators of the
Clifford group as well.

The last Clifford generator matrix is a two-qubit gate (a $4 \times 4$ unitary matrix) which
also represents a \emph{controlled} operation. That is, based on the
$\ket{1}$ component of the \emph{control} qubit, it applies a single-qubit
gate (in this case, Pauli $X$) to the \emph{target} qubit.
In fact,
both CNOT and $X$ are also fundamental gates in classical reversible
logic as well, where $X$ is also the Boolean $NOT$ gate on classical bits.
That is why the gate is called CNOT, for ``controlled-NOT.'' Its inclusion
in the generating set for $\mathcal{C}_2$ means that it can be applied
in either direction: with control on qubit 1 and target on qubit 2 or
vice versa. CNOT is defined below.
%
\begin{equation}
CNOT = 
 \left[
  \begin{array}{cccc}
    1 & 0 & 0 & 0 \\
    0 & 1 & 0 & 0 \\
    0 & 0 & 0 & 1 \\
    0 & 0 & 1 & 0
  \end{array} \right]
\end{equation}
%
Likewise, the general Clifford group on $n$ qubits $\mathcal{C}_n$
can be generated from the same set
as $\mathcal{C}_2$, with gates understood to apply to any of the $n$ qubits.
With three CNOTs, we can implement the $SWAP$ gate which exchanges the
states of two qubits.
%
\begin{equation}
SWAP = 
 \left[
  \begin{array}{cccc}
    1 & 0 & 0 & 0 \\
    0 & 0 & 1 & 0 \\
    0 & 1 & 0 & 0 \\
    0 & 0 & 0 & 1
  \end{array} \right]
\end{equation}
%
The gate CNOT has historical importance in quantum computing partly
due to its use in many
early quantum gate decompositions and its ability to
be performed fault-tolerantly in many physical technologies. It will be our
primary two-qubit gate.
Along with arbitrary single-qubit gates, it is universal for quantum computation \cite{Barenco1995a}.
Therefore, we will give this basis a special name:
%
\begin{equation}
\mathcal{Q} = \{ U(2) \cup CNOT \}
\end{equation}

%%%%%%%%%%%%%%%%%%%%%%%%%%%%%%%%%%%%%%%%%%%%%%%%%%%%%%%%%%%%%%%%%%%%%%%%%%%%%%
\subsection{Controlled Gates}
\label{subsec:controlled}

The principle of a controlled gate can be generalized to multiple
controls using the ``meta-operator'' notation from \cite{Kitaev2002}.
By $\Lambda^n(U)$, we mean an $(n+1)$-qubit gate ($2^{n+1} \times 2^{n+1}$
unitary matrix) with $n$ control qubits and a single-qubit target gate
$U \in U(2)$. An important multiply-controlled gate, which is universal
for classical reversible circuits, is the Toffoli gate, or controlled-controlled-$NOT$.
%
\begin{equation}
\text{Toffoli} = \Lambda^2(X) = 
 \left[
  \begin{array}{cccccccc}
    1 & 0 & 0 & 0 & 0 & 0 & 0 & 0 \\
    0 & 1 & 0 & 0 & 0 & 0 & 0 & 0 \\
    0 & 0 & 1 & 0 & 0 & 0 & 0 & 0 \\
    0 & 0 & 0 & 1 & 0 & 0 & 0 & 0 \\
    0 & 0 & 0 & 0 & 1 & 0 & 0 & 0 \\
    0 & 0 & 0 & 0 & 0 & 1 & 0 & 0 \\
    0 & 0 & 0 & 0 & 0 & 0 & 0 & 1 \\
    0 & 0 & 0 & 0 & 0 & 0 & 1 & 0
  \end{array} \right]
\end{equation}
%
As seen above, multiply-controlled single-qubit gates $\Lambda^n(U)$ have a
special, sparser structure than general $n$-qubit gates in $U(2^n)$.
These play a special role in many multi-qubit decompositions, about which
we will say more in Chapter \ref{chap:qcompile}.

There is also a special case of a ``targetless'' controlled single-qubit
gate which simply rotates the $\ket{1}$ component of a single-qubit state.
%
\begin{equation}
\Lambda(e^{i\phi}) = 
 \left[
  \begin{array}{cc}
    1 & 0 \\
    0 & e^{i\phi} \\
  \end{array} \right]
\end{equation}
%
This gate is a key tool and simplification for single-qubit compiling.

%%%%%%%%%%%%%%%%%%%%%%%%%%%%%%%%%%%%%%%%%%%%%%%%%%%%%%%%%%%%%%%%%%%%%%%%%%%%%%
\subsection{Single-Qubit Compiling}
\label{subsec:qcompile-single}

A seemingly simpler task than general compiling is single-qubit compiling.
This will illustrate the basic principles of quantum compiling and the
structure that we will exploit later to choose an effective basis. Moreover,
it will reveal a general relationship between many of the single-qubit
gates that we have already introduced.

First, we review how to decompose a general $U \in U(2)$ into three single-qubit
rotations about the Bloch sphere $x$-axis and $z$-axis, the so-called
Euler angle decomposition \cite{Nielsen2000}. This gives rise to a factor of $3$
which commonly appears in resource calculations in the literature.
%
\begin{equation}
U = e^{i\delta}R_Z(\gamma)R_X(\beta)R_Z(\alpha)
\end{equation}
%
The gate $R_Z(\phi)$ represents a rotation about the Bloch sphere $z$-axis,
of which the Pauli $Z$ gate is a special case of a $\pi$ rotation. In fact,
it is the same as the controlled-phase gate we introduced in the previous section,
up to a global phase.
%
\begin{equation}
R_Z(\phi) =
\left[
  \begin{array}{cc}
    e^{-i\phi/2} & 0 \\
    0 & e^{i\phi/2} \\
  \end{array} \right]
=
e^{-i\phi / 2} \left[
  \begin{array}{cc}
    1 & 0 \\
    0 & e^{i\phi} \\
  \end{array} \right]
= e^{-i\phi / 2} \Lambda(e^{i\phi})
\end{equation}
%
We can now state the relationship between $S$ and $Z$, as well as introduce
an important new gate $T$ which is the square root of $S$ up to a phase. All three
are rotations about the Bloch $z$-axis by power-of-two fractions of $\pi$.
%
\begin{equation}
Z = R_Z(\pi) =
\left[
  \begin{array}{cc}
    1 & 0 \\
    0 & -1 \\
  \end{array} \right]
\qquad
S = R_Z(\tfrac{\pi}{2}) =
\left[
  \begin{array}{cc}
    1 & 0 \\
    0 & i \\
  \end{array} \right]
\qquad
T = R_Z(\tfrac{\pi}{4}) =
\left[
  \begin{array}{cc}
    1 & 0 \\
    0 & e^{i\pi / 4} \\
  \end{array} \right]
\end{equation}
%
Likewise, the gate $R_X(\phi)$ represents a rotation about the Bloch sphere $x$-axis,
of which the Pauli $X$ gate is a special case of a $\pi$ rotation.
%
\begin{equation}
R_X(\phi) =
\left[
  \begin{array}{cc}
    \cos \tfrac{\phi}{2} & -i \sin \tfrac{\phi}{2} \\
    -i \sin \tfrac{\phi}{2} & \cos \tfrac{\phi}{2} \\
  \end{array} \right]
\end{equation}
%
Similar decompositions can be given in terms of $R_X$ and $R_Y$, or in
terms of $R_Y$ and $R_Z$. Solving for the angles $\{ \alpha, \beta, \gamma, \delta \}$
involves writing four equations in four variables, which can be found in
the standard textbook \cite{Nielsen2000}. We will not say
more about their solution here, except that we can implement the
global phase shift $e^{i\delta}$ using the identities below, which are
adapted from \cite{Kitaev2002}.
%
\begin{eqnarray}
e^{i\delta} & = & R_Z(\phi)\cdot X \cdot R_Z(\phi) \cdot X \\
X & = & R_X(\pi) \\
Z & = & R_Z(\pi) \\
R_X(\phi) & = & H \cdot R_Z(\phi) \cdot H
\end{eqnarray}
%
It now seems that a reasonable basis for single-qubit compiling are
arbitrary $R_Z(\phi)$ and $R_X(\phi)$ rotations, along with $H$.
However, in practice our classical control can only implement
rotations with finite precision. How can we measure this, or any
other, precision arising from approximation?

%%%%%%%%%%%%%%%%%%%%%%%%%%%%%%%%%%%%%%%%%%%%%%%%%%%%%%%%%%%%%%%%%%%%%%%%%%%%%%
\subsection{Distance Metrics}
\label{subsec:distance}

Each gate from our basis is a unitary matrix of bounded dimension, and the action of an entire
$n$-qubit compiled circuit $\tilde{C}$
is also a $2^n \times 2^n$ unitary matrix. This matrix can be formed
by the product of $2^n \times 2^n$ matrices $G_i$ which are a tensor
product of gate matrices from our basis operating on individual qubits.
Our desired target matrix $C$ is itself
a matrix from $U(2^n)$, and therefore we will need a distance metric
that operates on matrices (specifically, the difference of matrices).
%
\begin{equation}
\tilde{C} = G_{D}G_{D-1}\cdots G_{2} G_{1}
\end{equation}
%
One distance metric used in theoretical literature
is the operator norm of a matrix $M$,
which is defined as the maximum amount $M$ scales the vector norm
of all unit-length vectors. This is sometimes also called the
infinity-norm, or supremum-norm (sup-norm).
%
\begin{equation}
\| M \|_{\infty} = \max_{\| \ket{v} \| = 1} \| M \ket{v} \|
\label{eqn:op-norm}
\end{equation}
%
However, this is not an operational definition.
Moreover, we often wish to neglect a global phase in a unitary matrix,
which is not measurable in quantum physics. To measure phase-independent
distance between two unitary matrices, we can use the following
distance measure due to Fowler \cite{Fowler2011}.
%
\begin{equation}
dist(U, V) = \| U - V\| \equiv \sqrt{\frac{2^n - |tr(U^{\dag}V)|}{2^n}}
\end{equation}
%
Now can quantify the quality of our approximations through an
error $\epsilon$.
%
\begin{equation}
\| C - \tilde{C}\| = \| C -  G_D\cdots G_1 \| < \epsilon
\end{equation}
%
Often $\epsilon$ will be small, and we will upper bound it by some
power of $\frac{1}{2}$. Therefore, we define a
new parameter $\nu$ which is the number of bits needed to encode the exponent
of this increasingly small fraction.
%
\begin{equation}
\epsilon = \frac{1}{2^\nu} \qquad
\nu = \log(1/\epsilon)
\end{equation}
%
It is natural to suppose that compiling better approximations requires
more resources, and these resources are expressed as functions of
these
parameters, increasing with decreasing $\epsilon$ and increasing $\nu$.
In fact, often the efficiency and the
capabilities of a quantum
compiler depend on its basis. Therefore, we conclude this section by
discussing circuit bases.

%%%%%%%%%%%%%%%%%%%%%%%%%%%%%%%%%%%%%%%%%%%%%%%%%%%%%%%%%%%%%%%%%%%%%%%%%%%%%%
\subsection{Circuit Bases}
\label{subsec:qcompile-bases}

\begin{definition}{\textbf{Circuit basis.}}
A basis for a quantum circuit (family) is a universal set of
bounded-qubit gates (usually operating on no more than 3 qubits each).
We call a basis \emph{finite} if it contains a finite
number of gates; that is, it contains discrete gates and not an infinite
continuum of gates. We call a basis \emph{fixed} if its members are independent
of the number of qubits in the input circuit (its width).
\end{definition}

For fault-tolerant quantum computing, we are interested in compiling
circuits to a fixed, finite basis. What does it mean for a fixed, finite
basis to be universal for an infinite group like $SU(2^n)$?

\begin{definition}{\textbf{Universal approximation.}}
We call a fixed, finite set of gates $\mathcal{G}$ \emph{universal} for
a group $U(2^n)$ iff for every desired target $C \in U(2^n)$ and
desired error $\epsilon$, we can return a
sequence of gates $(g_1,g_2,\ldots,g_S)$ from $\mathcal{G}$ where
%
\begin{equation}
\| C - g_1 g_2 \cdots g_S \| \le \epsilon
\end{equation}
%
\end{definition}

This defines whether a gateset is a basis, or whether universal approximation
is even possible (non-constructively). We will see in
Chapter \ref{chap:qcompile} that
quantum compilers are concerned with constructive approaches to
\emph{efficiently} return such a compiled sequence $\tilde{C} = \prod g_i$.
Examples of efficiency metrics are the number of returned gates $S$ or
the running time of the quantum compiling algorithm (which itself is
primarily classical).

What gatesets are known to be fixed, finite, and universal, and therefore
suitable bases for quantum compilation? We state without proof that
the following basis satisfies these properties, and we will use it later as
our circuit basis in Chapters \ref{chap:factor-polylog} and
\ref{chap:factor-sublog}.
%
\begin{equation}
\mathcal{G} = \{X, Z, H, \text{Toffoli}, \text{CNOT}\}
\end{equation}
%
It is important to note that the only non-Clifford gate in the above basis
is Toffoli.
The Clifford group $\mathcal{C}_n$ by itself is \emph{not}
universal.
In fact, it is provable
that \emph{any} universal gateset must possess at least one
non-Clifford gate \cite{Zeng2011}.
Two popular choices for the non-Clifford gate in a basis are the $T = R_Z(\pi/4)$
gate and the $\text{Toffoli} = \Lambda^2(X)$ gate. Since these are not ``natively''
supported (non-transversal) in many codes, they must often be implemented
probabilistically using only Clifford operations and $MeasureZ$, usually
by way of a so-called ``magic'' state.
Therefore, many quantum compilers
use the Clifford+$T$ basis ($\mathcal{C}_2 \cup \{ T \}$)
or the Clifford+Toffoli basis ($\mathcal{C}_2 \cup \{ \text{Toffoli} \}$),
and measure
the non-Clifford gate as the most expensive resource. It is an area of
active research
whether $T$ or Toffoli is more efficient to implement
\cite{Jones2013a,Eastin2012}.

For single-qubit compilation, the $\{H,T, T^{\dagger}\}$ gateset is universal and
plays an important role in the literature. Other compilers may add
the Clifford gates $S$ and $S^{\dagger}$ to the bases above,
but this does not change their universality nor its asymptotic efficiency for
compiling. From our discussion in this section
on quantum gates and circuit bases, we are now prepared to add some additional
circuit constraints to form quantum architectures.

%%%%%%%%%%%%%%%%%%%%%%%%%%%%%%%%%%%%%%%%%%%%%%%%%%%%%%%%%%%%%%%%%%%%%%%%%%%%%%%
\section{Constant-depth Communication}
\label{sec:intro-cdc}

Communication, namely the \emph{moving} and \emph{copying} of quantum information, in nearest-neighbor quantum architectures is challenging.
In this section we quote known results for teleportation and
fanout in constant depth while also contributing a novel construction
for unfanout.

The first challenge of moving quantum information from one site to another over
arbitrarily long distances can be addressed by using
%A related problem is how to teleport a qubit an arbitrary distance.
% in an
%architecture through ancillae prepared in some initial state.
the constant-depth teleportation circuit
shown in Figure \ref{fig:cdt} due to Rosenbaum \cite{Rosenbaum2012}, illustrated using standard quantum circuit
notation \cite{Nielsen2000}. This requires the circuit resources shown in
Table \ref{tab:cd-resources}. The depth includes a layer of $H$ gates; a layer of CNOTs; an interleaved layer of Bell basis measurements; and two layers of
Pauli corrections ($X$ and $Z$ for each qubit), occurring concurrently with
resetting the $\ket{j}$ and $\ket{k}$ qubits back to $\ket{0}$.
These correction layers are not shown in the circuit.

\begin{figure*}[tb!]
\begin{center}
\begin{displaymath}
%\begin{array}{ccc}
\Qcircuit @C=1em @R=1em {
\lstick{\ket{\psi}}	& \qw      & \qw      & \qw & \qw & \qw & \qw & \qw                                          & \qw & \qw & \multimeasureD{1}{\mbox{Bell}} & \cw & \rstick{j_1} \\
\lstick{\ket{0}}    & \gate{H} & \ctrl{1} & \qw & \qw & \qw & \qw & \qw                                          & \qw & \qw & \ghost{\mbox{Bell}}            & \cw & \rstick{k_1} \\
\lstick{\ket{0}}    & \qw      & \targfix & \qw & \qw & \qw & \qw & \qw_{Z^{j_1}X^{k_1}\ket{\psi}}               & \qw & \qw & \multimeasureD{1}{\mbox{Bell}} & \cw & \rstick{j_2} \\
\lstick{\ket{0}}    & \gate{H} & \ctrl{1} & \qw & \qw & \qw & \qw & \qw                                          & \qw & \qw & \ghost{\mbox{Bell}}            & \cw & \rstick{k_2} \\
\lstick{\ket{0}}    & \qw      & \targfix & \qw & \qw & \qw & \qw & \qw_{Z^{j_2}Z^{j_1}X^{k_2}X^{k_1}\ket{\psi}} & \qw & \qw & \multimeasureD{1}{\mbox{Bell}} & \cw & \rstick{j_3} \\
\lstick{\ket{0}}    & \gate{H} & \ctrl{1} & \qw & \qw & \qw & \qw & \qw                                          & \qw & \qw & \ghost{\mbox{Bell}}            & \cw & \rstick{k_3} \\
\lstick{\ket{0}}    & \qw      & \targfix & \qw & \qw & \qw & \qw & \qw & \qw_{Z^{j_1}Z^{j_2}Z^{j_3}X^{k_3}X^{k_2}X^{k_1}\ket{\psi}} & \qw & \qw              & \qw & \qw \\
}
\end{displaymath}
\centerline{}
\caption{Constant-depth circuit based on \protect{\cite{Broadbent2007,Browne2009}} for teleportation over $n=5$ qubits \protect{\cite{Rosenbaum2012}}.}
\label{fig:cdt}
\end{center}\end{figure*}

Although general cloning is
impossible \cite{Nielsen2000}, the second challenge of copying information can be addressed by performing an unbounded quantum
fanout (or just fanout) operation:
$\ket{x,y_1,\ldots,y_n} \rightarrow \ket{x,y_1\oplus x, \ldots, y_n\oplus x}$.
This is used in our arithmetic circuits when
a single qubit needs to control (be entangled with) a large quantum register
(called a \emph{fanout rail}). The result of a fanout is an $n$-qubit
\emph{cat state} stored in the fanout rail, as shown in the following equation.

\begin{equation}
\normtwo (\ket{0}^{\otimes n} + \ket{1}^{\otimes n})
\label{eqn:fanned-out}
\end{equation}

To implement the fanout from the basis $\mathcal{G}$,
we employ a constant-depth circuit due to insight from
measurement-based quantum computing \cite{Raussendorf2003}
that first relies on the creation of an
$n$-qubit cat state \cite{Browne2009}. This method was communicated to
us by Harrow and Fowler \cite{Harrow2012}.
This circuit requires $O(1)$-depth, $O(n)$-size, and $O(n)$-width. Approximately
two-thirds of the ancillae are reusable and can be reset to $\ket{0}$ after
being measured. Numerical upper bounds are given in Table \ref{tab:cd-resources}.
The constant-depth fanout circuit is shown in Figure \ref{fig:cdf}
for the case of fanning out a given single-qubit state
$\ket{\psi} = \alpha\ket{0} + \beta\ket{1}$ to four qubits.

\begin{figure}[tb!]
\begin{center}
\begin{displaymath}
%& \qquad \qquad \qquad &
\Qcircuit @C=1em @R=1em {
\lstick{\ket{\psi}}	& \qw      & \qw      & \qw & \qw & \qw & \multimeasureD{1}{\mbox{Bell}'} & \cw & \rstick{j_1} \\
\lstick{\ket{0}}    & \gate{H} & \ctrl{1} & \qw & \qw      & \qw & \ghost{\mbox{Bell}'}            & \cw & \rstick{k_1} \\
\lstick{\ket{0}_1}    & \qw      & \targfix & \qw & \ctrl{1} & \qw & \qw      & \qw & \rstick{Z^{j_1}X^{k_1}\ket{\ell}_1}\\
\lstick{\ket{0}}	& \qw      & \qw      & \qw & \targfix & \qw & \multimeasureD{1}{\mbox{Bell}} & \cw & \rstick{j_2} \\
\lstick{\ket{0}}    & \gate{H} & \ctrl{1} & \qw & \qw      & \qw & \ghost{\mbox{Bell}}           & \cw & \rstick{k_2} \\
\lstick{\ket{0}_2}    & \qw      & \targfix & \qw & \ctrl{1} & \qw & \qw      & \qw & \rstick{Z^{j_2}X^{k_2}X^{k_1}\ket{\ell}_2}\\
\lstick{\ket{0}}	& \qw      & \qw      & \qw & \targfix & \qw & \multimeasureD{1}{\mbox{Bell}} & \cw & \rstick{j_3} \\
\lstick{\ket{0}}    & \gate{H} & \ctrl{1} & \qw & \qw      & \qw & \ghost{\mbox{Bell}}           & \cw & \rstick{k_3} \\
\lstick{\ket{0}_3}    & \qw      & \targfix & \qw & \ctrl{1} & \qw & \qw      & \qw & \rstick{Z^{j_3}X^{k_3}X^{k_2} X^{k_1}\ket{\ell}_3}\\
\lstick{\ket{0}_4}	& \qw      & \qw      & \qw & \targfix & \qw & \qw      & \qw & \rstick{X^{k_3}X^{k_2} X^{k_1}\ket{\ell}_4}\\
}
%& & \\
%(a) & & (b)
%\end{array}
\end{displaymath}
\centerline{}
\caption{Constant-depth circuits based on \protect{\cite{Broadbent2007,Browne2009}} for fanout \protect{\cite{Harrow2012}} of one qubit to $n=4$ entangled copies.}
\label{fig:cdf}
\end{center}
\end{figure}

The technique works by creating multiple small
cat states of a fixed size (in this case, three qubits), linking them
together into a larger cat state of unbounded size with Bell basis measurements,
and finally entangling them with the source qubit to be fanned out.
The qubits marked $\ket{\ell}$ are
entangled into the larger fanned out state given in Equation \ref{eqn:cat4}.
The Pauli corrections from the cat state creation are denoted by
$X^{k_2}$, $X^{k_3}$, $Z^{j_2}$ and $Z^{j_3}$ on qubits ending in
states $\ket{\ell}_1$, $\ket{\ell}_2$,
$\ket{\ell}_3$, and $\ket{\ell}_4$. The Pauli corrections
$X^{k_1}$ and $Z^{j_1}$ are from the Bell basis measurement
entangling the cat state with the source qubit (denoted $\text{Bell}'$).
\begin{equation}
Z_1^{j_1}X_1^{k_1}Z_2^{j_2}X_2^{k_2}X_2^{k_1}Z_{3}^{j_3}X_{3}^{k_3}X_{3}^{k_2}X_{3}^{k_1}X_{4}^{k_3}X_{4}^{k_2}X_{4}^{k_1}
\left(\alpha \ket{0}_1\ket{0}_2\ket{0}_3\ket{0}_4 + \beta \ket{1}_1\ket{1}_2\ket{1}_3\ket{1}_4 \right)
\label{eqn:cat4}
\end{equation}
%
The operators $X^k_i$ and $Z^j_{h}$ denote Pauli $X$ and $Z$ operators
on qubits $i$ and $h$, controlled by classical bits $k$ and $j$,
respectively. These corrections are enacted by the classical controller based on
the Bell measurement outcomes (not depicted).
Note the cascading nature of these corrections.
There can be up to
$n-1$ of these $X$ and $Z$
corrections on the same qubit, which can be simplified by the classical
controller to a single $X$ and $Z$ operation and then applied with a circuit of
depth 2 and size 2. Also, given the symmetric nature of the cat state, there
is an alternative set of Pauli corrections which would give the same state and
is of equal size to the corrections given above.

Reversing the fanout is an operation called \emph{unfanout}. Unfanout
takes as input 
the $n$-qubit cat state from Equation \ref{eqn:fanned-out} which is the result of a fanout.
The output of unfanout, after Pauli corrections, is the product state
consisting of all $\ket{0}$'s except for a single target qubit $\alpha\ket{0} + \beta\ket{1}$, which is in the
same state as the original source qubit of the fanout.

In the model of \cite{Hoyer2002}, the fanout and unfanout were identical, elementary
operations. In Figure \ref{fig:cdf}, the given \textsf{CCNTC} fanout circuit is not its own self-inverse
due to
the one-way nature of the measurement and the assumption that ancillary qubits begin in the
$\ket{0}$ state. Therefore, one major contribution of this current work is
a \textsf{CCNTC} circuit which performs unfanout, or the inverse of the fanout circuit in
Figure \ref{fig:cdf}. The relationship of our fanout and unfanout constructions and their
effect on intermediate quantum states is shown below.
\begin{equation}
(\alpha\ket{0} + \beta\ket{1})\otimes\ket{0}^{\otimes n-1} \rightarrow^{fanout}
\alpha\ket{0}^{\otimes n} + \beta\ket{1}^{\otimes n} \rightarrow^{unfanout}
(\alpha\ket{0} + \beta\ket{1})\otimes\ket{0}^{\otimes n-1}
\end{equation}
We give a concrete example of our unfanout circuit in Figure \ref{fig:cdu}
for $n=7$ which executes in constant depth on \textsf{2D CCNTC}.
Note that the state in Equation \ref{eqn:fanned-out}
is completely symmetric in that all qubits are
equivalent entangled copies of each other. Therefore, the asymmetry 
of the final target qubit is entirely determined by the unfanout circuit,
which in this case selects the bottom qubit in the figure to be the target.

\begin{figure*}[tb!]
\begin{center}
\begin{displaymath}
\Qcircuit @C=1em @R=1em {
& \lstick{\ket{\ell}}	& \qw & \gate{H} & \qw & \ctrl{1} & \qw & \qw      & \qw &  \measureD{Z} & \cw & \rstick{j_1} & \\
& \lstick{\ket{\ell}}	& \qw & \gate{H} & \qw & \targfix & \qw & \ctrl{1} & \qw & \measureD{Z} & \cw & \rstick{j_2} & \\
& \lstick{\ket{\ell}}	& \qw & \gate{H} & \qw & \ctrl{1} & \qw & \targfix & \qw & \measureD{Z} & \cw & \rstick{j_3} & \\
& \lstick{\ket{\ell}}	& \qw & \gate{H} & \qw & \targfix & \qw & \ctrl{1} & \qw & \measureD{Z} & \cw & \rstick{j_4} & \\
& \lstick{\ket{\ell}}	& \qw & \gate{H} & \qw & \ctrl{1} & \qw & \targfix & \qw & \measureD{Z} & \cw & \rstick{j_5} & \\
& \lstick{\ket{\ell}}	& \qw & \gate{H} & \qw & \targfix & \qw & \ctrl{1} & \qw & \measureD{Z} & \cw & \rstick{j_6} & \\
& \lstick{\ket{\ell}}	& \qw & \gate{H} & \qw & \qw      & \qw & \targfix & \qw & \gate{H} & \qw & \rstick{Z^{j_2 \oplus j_4}(\alpha\ket{0} + \beta\ket{1})}
}
\end{displaymath}
\centerline{}
\caption[A novel, constant-depth circuit for unbounded quantum unfanout on
\textsf{CCNTC}]{A novel, constant-depth circuit for unbounded quantum unfanout on
\textsf{CCNTC}, from the $7$-qubit entangled state $\alpha\ket{0}^{\otimes 7} + \beta\ket{1}^{\otimes 7}$ to the
target product state $(\alpha\ket{0} + \beta\ket{1})\otimes\ket{0}^{\otimes 6}$.}
\label{fig:cdu}
\end{center}\end{figure*}

The initial fanned out state lives in a $2$-dimensional subspace. The
round of Hadamard gates increases its dimension to $2^n$, and the two
interleaved layers of CNOTs in a sense ``disentangle'' the qubits from
one another, up to a Pauli $Z$ correction. This correction, on the
final target qubit, is controlled by the parity of the classical measurements
on every ``even'' qubit ($j_2$ and $j_4$ in the figure), excluding the 
next-to-last qubit ($j_6$ in the figure). Each measurement projects the state of the
target qubit
into a subspace with half the dimension, so $n-1$ measurements project
the target qubit into a final $2$-dimensional subspace, which is the
qubit $\alpha\ket{0} + \beta\ket{1}$.

Although the circuit shown works for odd $n$, we can easily take into
account even $n$ with an initial CNOT to ``uncopy'' one qubit from its
neighbors. The unfanout circuit in Figure \ref{fig:cdu} is the
functional inverse of the fanout circuit in 
Figure \ref{fig:cdf}, which by itself only requires a \textsc{1D} layout.
However, it relies on the fanned-out qubits
being teleported back into adjacent positions (in constant depth),
which is only possible in an overall \textsc{2D} layout.
This layout is not shown but can easily be constructed.
The target qubit of an unfanout is usually chosen to be in the same location
as the source qubit of the previous, corresponding fanout. 
The resources for unfanout are given in
Table \ref{tab:cd-resources}.

% From Notebook #16, p. 212
% From Notebook #16, p. 66
\begin{table}
\begin{displaymath}
\begin{tabular}{|c|c|c|c|}
\hline
\text{Circuit Name} & \text{Depth} & \text{Size} & \text{Width}\\
\hline
\text{Teleportation from Figure \ref{fig:cdt}} & $7$ & $3n + 4$ & $n+1$\\
\hline
\text{Fanout from Figure \ref{fig:cdf}} & $9$ & $10n - 9$ & $3n-1$ \\
\hline
\text{Unfanout from Figure \ref{fig:cdu}} & $ 6 $ & $ 3n+2 $ & $ n$ \\
\hline
\end{tabular}
\end{displaymath}
\centerline{}
\caption{Circuit resources for teleportation, fanout, and unfanout.}
\label{tab:cd-resources}
\end{table}

From an experimental perspective, it is physically efficient to create
a cat state in trapped ions using the M{\o}lmer-S{\o}rensen gate
\cite{Sorensen2000}\cite{Benhelm2008}. However, the fanout circuit for
the \textsf{CCNTCM} model would still be useful for other technologies, such
as superconducting qubits on a \textsc{2D} lattice.


\section{Quantum Architecture}
\label{sec:fpl-bg}

Quantum architecture is the design of physical qubit layouts
and their allowed interactions to execute
quantum algorithms efficiently in time, space, and other
resources.
In this dissertation, we focus on designing a realistic nearest-neighbor circuit for running
Shor's factoring algorithm on two-dimensional
architectural models of a physical quantum device with nearest-neighbor
interactions.

%%%%%%%%%%%%%%%%%%%%%%%%%%%%%%%%%%%%%%%%%%%%%%%%%%%%%%%%%%%%%%%%%%%%%%%%%%%%%%%
\subsection{Architectural Models and Circuit Resources}
\label{subsec:models}

Following Van Meter and Itoh \cite{VanMeter2005},
we distinguish between a model and an architectural implementation as follows.
A \emph{model} is a set of constraints and rules for the placement and
interaction of qubits.
An \emph{architecture} (or interchangeably, an \emph{implementation} 
or a \emph{circuit}) is a particular
spatial layout of qubits (as a graph of vertices) and allowed interactions (edges between the vertices),
following the constraints of a given model. In this section, we describe
several models that try to incorporate resources of physical interest from
experimental work. We also introduce a new model,
\textsf{2D CCNTCM}, which we will use to analyze our current circuit.

The most general model is called Abstract Concurrent (\textsf{AC})
and allows arbitrary, long-range interactions between any qubits and concurrent
operation of quantum gates.
This corresponds to a complete graph with an edge between every pair of nodes.
It is the model assumed in most quantum algorithms.

A more specialized model restricts interactions to nearest-neighbor, two-qubit,
concurrent gates (\textsf{NTC}) in a regular one-dimensional chain (\textsf{1D NTC}),
which is sometimes called linear nearest-neighbor (\textsf{LNN}).
This corresponds to a line graph. This is a more realistic model than
\textsf{AC}, but correspondingly, circuits in this model may incur greater
resource overheads.

To relieve movement congestion,
we can consider a two-dimensional regular grid
(2D NTC), where each
qubit has four planar neighbors, and 
there is an extra degree of freedom over the 1D model
in which to move data.
In this paper, we extend the \textsf{2D NTC} model in three ways.
The first two extensions are described in Section \ref{subsec:2dccntc},
and the third extension is described in Section \ref{subsec:2dccntcm}.

\subsection{\textsf{2D CCNTC}: Two-Dimensional Nearest-Neighbor Two-Qubit Concurrent Gates with Classical Controller}
\label{subsec:2dccntc}

The first extension allows arbitrary planar graphs
with bounded degree, rather than a regular square lattice.
Namely, we assume qubits lie in a plane and edges are not allowed to intersect.
All qubits are accessible from above
or below by control and measurement apparatus.
Whereas 2D NTC conventionally assumes each qubit
has four neighbors, we consider up to six neighbors in a roughly hexagonal
layout. The edge length in this model is no more than twice the edge length
in a regular 2D NTC lattice. The second extension is the realistic assumption
that classical control (CC) can
access every qubit in parallel, and we do not count these classical
resources in our implementation since they are polynomially bounded. The
classical controllers
correspond to fast digital computers which are
available in actual experiments and are necessary for constant-depth
communication in the next section.

We call an AC or NTC model augmented by these two extensions
\textsf{CCAC} and \textsf{CCNTC}, respectively. Before we describe the
third extension, let us formalize our model for \textsf{2D CCNTC}, with definitions that are (asymptotically) equivalent to those in 
\cite{Rosenbaum2012}.

\begin{definition}
A \textsf{2D CCNTC} architecture consists of

\begin{itemize}
\item a quantum computer $QC$ which is represented by a planar graph $(V,E)$. A
node $v \in V$ represents a qubit which is acted upon in a circuit, and an
undirected edge $(u,v) \in E$ represents 
an allowed two-qubit interaction between qubits $u,v \in V$. Each node has
degree at most $6$.
\item a universal gate set $\mathcal{G} = \{X, Z, H, T, T^{\dagger}, CNOT, MeasureZ\}$.

\item a deterministic machine (classical controller) $CC$ that applies a sequence
of concurrent gates in each of $D$ timesteps.
\item In timestep $i$, $CC$ applies a set of
gates $G_i = \{g_{i,j} \in \mathcal{G} \}$.
Each $g_{i,j}$ operates in one of the following two ways:
\begin{enumerate}
\item It is a single-qubit gate from $\mathcal{G}$ acting on a single qubit $v_{i,j} \in V$
\item
It is the gate CNOT from $\mathcal{G}$ acting on two qubits $v^{(1)}_{i,j}, v^{(2)}_{i,j} \in V$ where
$(v^{(1)}_{i,j}, v^{(2)}_{i,j}) \in E$
\end{enumerate}
All the $g_{i,j}$ can only operate on
disjoint qubits for a given timestep $i$. We define the support of $G_i$
as $V_i$, the set of all qubits acted upon during timestep $i$.

\begin{equation}
V_i = \bigcup_{j: g_{i,j} \in G_i} v_{i,j} \cup v^{(1)}_{i,j} \cup v^{(2)}_{i,j}
\end{equation}

\end{itemize}
\end{definition}

We can then define the three conventional circuit resources in this model.

\begin{description}
\item[circuit depth ($D$):] the number of concurrent timesteps.
\item[circuit size ($S$):] the total number of non-identity gates applied
from $\mathcal{G}$, equal to $\sum_{i=1}^D |G_i|$.
\item[circuit width ($W$):] the total number of qubits operated upon by
any gate, including inputs, outputs, and ancillae. It is equal to $| \bigcup_{i=1}^D V_i|$.
\end{description}

We observe that the following relationship holds between the circuit resources.
The circuit size is bounded above by
the product of circuit depth and circuit width, since in the worst case,
every qubit is acted upon by a gate for every timestep of a circuit.
The circuit depth is also bounded above by the size, since in the worst case,
every gate is executed serially without any concurrency.

\begin{equation}
D \le S \le D\cdot W
\label{eqn:depth-width}
\end{equation}

The set $\mathcal{G}$ includes measurement in the $Z$ basis, which is
actually not a unitary operation but which may be slower than unitary
operations in actual practice \cite{DiVincenzo2007}.
Therefore we count it in our resource
estimates.
All other gates
in $\mathcal{G}$ form a universal set of unitary
gates \cite{Kitaev2002}.
 In this paper we
will treat the operations in $\mathcal{G}$ as \emph{elementary gates}.
We can also define a Bell basis measurement using operations
from $\mathcal{G}$. A circuit performing this measurement is shown
in Figure \ref{fig:bell-measure} and has depth $4$,
size $4$, and width $2$.

\begin{figure*}[tb!]
\begin{center}
\begin{displaymath}
\begin{array}{ccc}
\Qcircuit @C=1em @R=1em {
& \qw & \multimeasureD{1}{\mbox{Bell}} & \cw & \rstick{j} \\
& \qw & \ghost{\mbox{Bell}}            & \cw & \rstick{k}
}
& \qquad \equiv \qquad &
\Qcircuit @C=1em @R=1em {
& \qw & \ctrl{1} & \qw & \gate{H} & \qw & \meter & \cw & \rstick{j} \\
& \qw & \targfix & \qw & \qw      & \qw & \meter & \cw & \rstick{k}
}
\end{array}
\end{displaymath}
\centerline{}
\caption{A circuit for measurement in the Bell state basis.}
\label{fig:bell-measure}
\end{center}\end{figure*}

The third extension to our model, and the most significant, is to consider
multiple disconnected planar graphs, each of which is a 2D CCNTC
architecture. This is described in more detail in the next section.

\subsection{Random Notes on 2D Architectures}

The main point should be shifted, in the intro, and maybe elucidated here,
that low-depth for factoring can be achieved by moving into 2D, but
moreover that the width and size blowups can also be reduced with a
hybrid 2D module.

The model is hybrid, because instead of a single, contiguous 2D lattice,
say with $W$ qubits and $S$ interactions between qubits (counting the
general case of all single- and two-qubit gates as interactions),
the computation is split among many communicating 2D lattices
(say $\overline{W}$ of them), each containing $W/\overline{W}$ qubits
each. In this new model, we argue that there are now a total number of
qubit interactions $S'$ across all modules, where $S' \le S$, and at
the same time we have introduced $\overline{S}$
new, \emph{long-range} interactions between modules. Based on experimental
data of create shared entanglement between separate ion traps (citation needed here),
these long-range interactions are more expensive than short-range interactions
within a module, and so we count them separately.

Currently, we impose no constraints on the connectivity of the modules, which exist
on a higher-level graph of modules. This is supported by current proposals
(cite MUSIQC here). However, it is interesting to see what effect such constraints
would have on a factoring architecture. We will return to this later.

With no such constraints (arbitrary connectivity of modules), we can now see the
effect of setting the size of each module. If we set modules
to be a single qubit in size ($\overline{W} = W$), we would get the architectural
model \textsf{AC}. If we were to set module size to include the entire lattice
($\overline{W} = 1$), then we would get the model described in the previous
section, \textsf{2D CCNTC}. This new model then, \textsf{2D CCNTCM}, is a hybrid
between \textsf{AC} and \textsf{2D CCNTC} with module size as a parameter.
\textsf{AC} and \textsf{2D CCNTC} represent the two extremes of long-range
interactions, where \textsf{AC} permits all long-range interactions and
\textsf{2D CCNTC} permits none.

There are now two more things to study and discuss. One is the optimal setting
of the module size $W/\overline{W}$, which in this paper we have set to
be $O(n)$. The second is the connectivity constraints of the modules.

For the first parameter, we have not rigorously shown that the optimal setting
of $W / \overline{W}$ is 

Now then, what is the effect of imposing constraints on module connectivity?

\subsection{\textsf{2D CCNTCM}: Two-Dimensional Nearest-Neighbor Two-Qubit Concurrent Gates with Classical Controller and Modules}
\label{subsec:2dccntcm}

A single, contiguous
2D lattice that contains an entire quantum architecture may be prohibitively large to manufacture. In practice,
scalable experiments will probably use many
smaller quantum computers that communicate by means of shared
entanglement \cite{Monroe2012}.
We call these individual machines \emph{modules}, each of
which is a self-contained \textsf{2D CCNTC} lattice. This should not be
confused with the word ``modular'' as in ``modular arithmetic'' or as
referring to the modulus $m$ that we are trying to factor.

We treat these modules
and teleportations between them as nodes and edges, respectively,
in a higher-level planar graph. The teleportations each transmit one qubit
from one module to another, from any location within the source module
to any location within the destination module, making use of the
omnipresent classical controller. The modules can be arbitrarily far
apart physically, but they have bounded-degree connectivity with other
modules, and their edges are planar (they cannot intersect).

A single module can be part of multiple teleportation operations in a single timestep, as long as they involve disjoint qubits within the module.
We justify this assumption in that it is
possible to establish entanglement between multiple
quantum computers
in parallel. We call this new model \textsf{2D CCNTCM},
and we argue that is captures the essential aspects of 2D architectures
without being overly sensitive to the exact geometry of the lattices involved.
A graphic depiction of three modules in \textsf{2D CCNTCM} is shown in
Figure \ref{fig:modules}. Each module contains within it a
\textsf{2D CCNTC} lattice. We can equivalently consider the omnipresent,
single
classical controller as a collection of multiple classical controllers, one
for each module or teleportation operation, which can inter-communicate
classically and share a clock.

\textsf{2D CCNTCM} is very similar to the model of a quantum multicomputer
of Van Meter
\cite{VanMeter2006} with the difference that we allow parallel teleportations.
It is also very similar to the distributed quantum computer \textsf{DQC} of
Beals et al. \cite{Beals2012} except that we require each module to have
$\Omega(n)$ qubits in order to enforce some locality.

\begin{figure}[btp!]
\begin{center}
\includegraphics[width=4in]{factor-polylog/figures/modules.pdf}
\end{center}
\caption{Three modules in the \textsf{2D CCNTCM} model}
\label{fig:modules}
\end{figure}

\begin{definition}
A \textsf{2D CCNTCM} architecture consists of

\begin{itemize}
\item a quantum computer $\overline{QC}$ which is represented by a planar graph $(\overline{V},\overline{E})$. A
node $\overline{v} \in \overline{V}$ represents a module, or a graph $(V,E)$
from a \textsf{2D CCNTC} architecture defined previously. It can have
unbounded degree.
An
undirected edge $(\overline{u},\overline{v}) \in \overline{E}$ represents an
allowed teleportation from any qubit in module $\overline{u}$ to
another qubit in module $\overline{v}$.
\item All modules are restricted to be linear in the number of their qubits:
$|V| = \Theta(n)$ for all $(V,E) \in \overline{V}$.
\item a universal gate set $\mathcal{G} = \{X, Z, H, T, T^{\dagger}, CNOT,
MeasureZ\}$
for the qubits \emph{within the same} modules which is the same as for \textsf{2D CCNTC},
and an additional operation $Teleport$ which only operates on qubits
\emph{in
different} modules.
\item a deterministic machine (classical controller) $\overline{CC}$ that applies a sequence
of concurrent gates in each of $D+\overline{D}$ timesteps.
This can be a separate classical controller
for every pair of modules.
\item In timestep $i$, $\overline{CC}$ applies
gates $G_i = \{g_{i,j} : g_{i,j} \in \mathcal{G} \lor g_{i,j} = Teleport \}$.
That is, there are two kinds of timesteps with respect to the kinds of gates
which operate within them.
\begin{enumerate}
\item In the first kind, gates are exclusively from $\mathcal{G}$, and
they operate within modules as described
for \textsf{2D CCNTC} above. We say there are $D$ such timesteps.
\item In the second kind, gates are exclusively $Teleport$ gates between two qubits $v^{(1)}_{i,j} \in \overline{v}_1$ and
$v^{(2)}_{i,j} \in \overline{v}_2$ for
(possibly non-distinct) modules $\overline{v}_1, \overline{v}_2 \in \overline{V}$.
Again, all such qubits much be distinct within a timestep.
We say there are $\overline{D}$ such timesteps.
\end{enumerate}

Again, we define the support of $G_i$
as $V_i$, the set of all qubits acted upon by any $g_{i.j}$, which
includes all the modules.
\begin{equation}
V_i = \bigcup_{j: g_{i,j} \in G_i} v_{i,j} \cup v^{(1)}_{i,j} \cup v^{(2)}_{i,j} 
\end{equation}

\end{itemize}
\end{definition}

We measure the efficiency of a circuit in this new module using not just
the three conventional circuit resources, but with three novel resources
based on modules.

\begin{description}

%, depicted in Figure \ref{fig:resources}:
\item[module depth ($\overline{D}$):] the depth of consecutive teleportations between modules.
\item[module size ($\overline{S}$):] the number of total qubits teleported between any two modules over all timesteps.
\item[module width ($\overline{W}$):] the number of modules whose qubits are
acted upon during any timestep.

\end{description}

%We can make an observation analogous to Equation \ref{eqn:depth-width} but
%for modules in Equation \ref{eqn:module-depth-width}.

%\begin{equation}
%\overline{D} \le \overline{S} \le \overline{D}\cdot \overline{W}
%\label{eqn:module-depth-width}
%\end{equation}

We note the following relationship between circuit width and
module width.

\begin{equation}
W = O(n\overline{W})
\label{eqn:module-width}
\end{equation}

This restriction imposes some locality on our model by constraining it to
nearest-neighbor gates within a linear-sized group of qubits, but allowing
it long-range teleportation to circumvent onerous geometric constraints.
Using the constant-depth communication in Section \ref{sec:cdc}, and for
the specific case of factoring, we
can simulate arbitrary connectivity between modules with only a polynomial
increase in the module size and a constant increase in module depth.

\subsection{Circuit Resource Comparisons}

Counting gates from $\mathcal{G}$ as having unit size and unit depth
is
an overestimate compared to the model in \cite{Kutin2006}, in which a
two-qubit gate has unit size and unit depth and
absorbs the depth and size of any adjacent single-qubit gates. We intend
for this more pessimistic estimate to reflect the practical difficulties
in compiling these gates using a non-Clifford gate in a fault-tolerant way,
such as the $T$ gate or the Toffoli gate
\cite{Fowler2011}.
%However, these difficulties may be mitigated by using
%Toffoli gates directly, which can be fault-tolerantly implemented using
%magic-state distillation according to recent works \cite{Eastin2012,Jones2013a}.

In both our resource counting method and that of \cite{Fowler2004,Kutin2006}, multiple gates acting on disjoint qubits
can occur in parallel during the same timestep. For each building block,
from modular addition to modular multiplication and finally to modular
exponentiation, we provide closed form equations upper-bounding the required circuit
resources as a function of $n$, the size of the modulus $m$ to be factored.
We will use the
term \emph{numerical upper bound} to distinguish these formulae from asymptotic
upper bounds.

It is possible to reduce the numerical constants with more detailed analysis,
which would be important for any physical implementation.
However, we have chosen instead to simplify the number of terms in the formulae
for the current work. We do not intend for these upper bounds to represent
the optimal or final work in this area.

The modular adder in Section \ref{sec:csa-mod-add} and its carry-save
subcomponents only occur within a single module, so we only give their
circuit resources in terms of circuit depth, circuit size, and circuit width. 
For the modular multiplier in
Section \ref{sec:csa-mod-mult} and the modular exponentiator in
Section \ref{sec:modexp}, we also give circuit resources in
terms of module depth, module size, and module width.